\documentclass[11pt]{article}
\usepackage[T1]{fontenc}
\usepackage[utf8]{inputenc}
\usepackage{lmodern,microtype,amsmath,amssymb,amsthm,mathtools}
\usepackage[margin=0.92in]{geometry}
\usepackage{booktabs,longtable,array,enumitem,xcolor,fancyhdr}
\usepackage{hyperref}
\hypersetup{colorlinks=true,linkcolor=blue!45!black,urlcolor=blue!50!black,citecolor=blue!45!black,pdftitle={Simultaneous quadratic coding for ZFC consistency},pdfauthor={AI-assisted research for Vladimir Reshetnikov}}
\setlength{\emergencystretch}{3em}
\setlength{\parindent}{0pt}
\setlength{\parskip}{0.35em}
\setlist{nosep,leftmargin=1.6em}
\newtheorem{theorem}{Theorem}[section]
\newtheorem{lemma}[theorem]{Lemma}
\newtheorem{proposition}[theorem]{Proposition}
\newtheorem{remark}[theorem]{Remark}
\newcommand{\N}{\mathbb N}
\newcommand{\Z}{\mathbb Z}
\newcommand{\Pos}{\mathbb Z_{>0}}
\newcommand{\Con}{\operatorname{Con}}
\newcommand{\pc}{\operatorname{pc}}
\newcommand{\pin}{\texttt{22e75dd6b87b1b199bb08c778d7cf075ae5aae08}}
% Replaced from final verified ledgers before release.
\newcommand{\NineDegree}{129,875,980}
\newcommand{\NineOps}{788}
\newcommand{\NineNatOps}{791}
\newcommand{\TenDegree}{237,526,776}
\newcommand{\TenOps}{936}

\pagestyle{fancy}\fancyhf{}
\fancyhead[L]{\small Simultaneous quadratic coding}
\fancyhead[R]{\small ZFC consistency}
\fancyfoot[C]{\thepage}\setlength{\headheight}{14pt}
\title{Simultaneous quadratic coding\\for ZFC consistency\\[0.5em]\large Short positional certificates, a multi-output carry test,\\and much lower-degree nine- and ten-variable equations}
\author{AI-assisted research prepared for Vladimir Reshetnikov\\ProveIt repository}
\date{23 October 2026\\\small Source snapshot: \texttt{22e75dd6b87b}}
\begin{document}
\maketitle
\begin{abstract}
The recent ProveIt constructions first combine a quadratic system into a
quartic and then assign exponentially separated positions to its monomials.
We replace both steps. A fixed 31-element positional certificate with largest
position 799 distinguishes every monomial of degree at most two. Twenty-two
separated target cells encode the quadratic residuals directly, and a single
binary no-carry condition tests every target at once. The resulting source
interface feeds the existing Jones and Sun arithmetic bridges.
The principal equations have nine positive variables, degree \NineDegree{}
and \NineOps{} binary operations, and ten unrestricted integer variables,
degree \TenDegree{} and \TenOps{} operations. Their predecessors had degrees
about $2.15\cdot10^{27}$ and $5.89\cdot10^{26}$, respectively.
An independent improvement gives twelve positive variables at degree 941
and 121 operations. We give complete polynomial specifications, the new
universal arithmetic proofs, exact finite certificates, fixed coefficient
payloads, and reproducible checks. The consistency equivalence is an ordinary
mathematical proof formalizable in ZFC, not a new Lean or Rocq certificate.
No global minimum or new universal unknown-count theorem is claimed.
\end{abstract}
\tableofcontents
\newpage

\section{Result and comparison}

Throughout, $\N=\{0,1,\ldots\}$ and $\Pos=\{1,2,\ldots\}$.
The domain is part of each assertion. A polynomial intended for positive
variables must not silently be interpreted over all integers.

\begin{theorem}[Main consistency presentations]\label{thm:main}
Let $F_9$ and $F_{10}$ be the fixed integer polynomials defined below and
bound to the complete coefficient instance in the archive. Then ZFC proves
\[
 \begin{aligned}
 \neg\exists\mathbf x\in\Pos^9\;F_9(\mathbf x)=0
 &\ \longleftrightarrow\ \Con(\mathrm{ZFC}),\\
 \neg\exists\mathbf z\in\Z^{10}\;F_{10}(\mathbf z)=0
 &\ \longleftrightarrow\ \Con(\mathrm{ZFC}).
 \end{aligned}
\]
Replacing the nine positive coordinates by nonnegative coordinates plus one
gives the supplied nine-natural-variable alternative.
\end{theorem}

The theorem is an equivalence, not a proof of either side. The internal
arithmetized consistency predicate is the same finite-refutation predicate
used by the calibrated K215 endpoint. No unconditional two-sided independence
claim follows merely from consistency, and none is made here.

\subsection{What is genuinely different}

There are three logically separate improvements.
\begin{enumerate}
\item Replace exponential Kronecker positions by a short, explicitly checked
additive positional certificate.
\item Encode all quadratic residuals simultaneously instead of squaring and
adding them into a quartic. The resulting source uses 136 weighted coefficient
records, rather than 402 quartic records.
\item Use a direct two-block carry source. Its positivity holds before the
final equation is solved, allowing substantially simpler square-root bounds
in the nine-variable finalizer. The enormous dense Jones index is no longer
needed.
\end{enumerate}
The first step uses classical Bose--Chowla ideas; the contribution is their
application through a new explicit source adapter. The second and third steps
are essential to the strongest concrete reductions. Merely substituting a
short list of exponents into an old theorem statement would not prove them.

\begin{center}\small
\begin{tabular}{@{}p{0.30\textwidth}rrrr@{}}
\toprule
Presentation & Domain & Degree & Old ops & New ops\\
\midrule
Main nine-variable equation & $\Pos^9$ & \NineDegree{} & 1500 & \NineOps{}\\
Ten-variable Sun instance & $\Z^{10}$ & \TenDegree{} & 1442 & \TenOps{}\\
Low-degree twelve-variable form & $\Pos^{12}$ & 941 & 124 & 121\\
Same twelve-variable system & $\Pos^{12}$ & 932 & 120 & 116\\
\bottomrule
\end{tabular}
\end{center}
For the twelve-variable comparison, the preceding lowest-degree scalar form
had degree 1022 and its system had maximum degree 1013. The new system has
three equations. The nine-variable predecessor in this comparison had degree
$92324\cdot5^{32}+296$, and the ten-integer predecessor had degree
$126496\cdot5^{31}+27064$.

Operations are live binary additions, subtractions and multiplications,
including multiplication by numerical constants. Fixed coefficient generation,
reuse of a computed value, and equality comparison in a system are free.
Exponentiation is expanded into these gates; no variable exponent is allowed.
This metric is not expanded monomial count, bit complexity, or coefficient
height. There is no single established meaning of ``simplest'' that makes all
these measures interchangeable.

\subsection{The latest seven packages}

We reviewed the corrected catalogue and packages 31--37 at commit
\texttt{22e75dd6b87b} (full pin in the source manifest).
The exact-source inventory records 83 Git blob checks for the initial review;
additional construction-specific pins accompany their source data.

\begin{longtable}{@{}p{0.11\textwidth}p{0.80\textwidth}@{}}
\toprule Package & Contribution retained or compared\\\midrule\endhead
31 & Sun's ten-integer instance; a 60-operation quadratic source system;
shared arithmetic in the twelve-variable norm; complete fixed tuple.\\
32 & Fifteen- and sixteen-integer variants using a different 32-variable
parent; cheaper positive norm evaluation.\\
33 & Sixteen-integer construction and careful restrictions on cofinal
square-base arguments.\\
34 & Nine-variable low-degree finalizer, signed constructions, leading-term
carrier corrections, and twelve-variable degree frontier.\\
35 & Alternative carrier cancellations and refactored positive norms.\\
36 & Smaller coherent compiler coefficient tuple, and rebinding of the
67-operation scalar and 59-operation systems.\\
37 & Faster nine-variable evaluation, fourteen-integer variants, and lower-cost
moderate-degree integer presentations.\\\bottomrule
\end{longtable}

The new constructions do not reduce the number of unknowns below the existing
nine-positive and ten-integer counts. The old 67-operation scalar equation and
59-operation systems remain useful and are not beaten on their own operation
measure. We retain package 31's complete tuple for a transparent like-for-like
source comparison; we do not claim the smallest coefficient radix.

\section{Fixed source, constants, and the verification boundary}

The common starting point is the inherited 31-positive-variable system of
22 quadratic equations. In its 60-operation coordinate version the variables
are
\[
\begin{gathered}
J,F,\alpha,z,f,h,i,o,j,s,w,\tau,\eta,\zeta,y,\gamma,\\
N,E,P,c,a,X,D,C_2,L,F_2,V,g,I,B_a,\Delta.
\end{gathered}
\]
The four fixed ports are denoted $\beta,K_0,M_C,M_F$; the last is already the
shifted high mask. Put
\[
\begin{gathered}
q=\beta J+1,\quad k=\eta+\zeta,\quad H=4a+3,\quad C=q-F-\alpha,\\
U=L+1,\quad Y=a-E,\quad Q=c-\eta,\quad T=P-k,\quad R=jc-V.
\end{gathered}
\]
Then the equations are
\[\begin{gathered}
U=q^2,\qquad N=qU\\
X=wq,\qquad Y=sN,\\
E=XY,\qquad Q=kY\\
T=EQ,\qquad \tau^2=PT+1,\\
D=X+ac+\gamma H,\qquad \Delta=a^2+H\\
C_2=c^2,\\
D^2=\Delta C_2+1,\qquad k=R+1+hE,\\
R=(q-F)(N-q)-CL+(M_C+qM_F)J,\\
(K_0+w)C=F+z\beta J,\\
F_2=f^2,\qquad V=of-c\\
g=y^2-V^2,\\
I=iC_2,\qquad B_a=I^2\\
B_a=\Delta(F_2-1),\\
y^2-B_ag=1.
\end{gathered}\]
Let $r_1,\ldots,r_{22}$ denote the left sides minus the right sides, in
this order. Replace each of the 31 positive coordinates by $n_i+1$ and
write $q_j(\mathbf n)=r_j(\mathbf n+\mathbf1)$. In particular,
$q_5(\mathbf0)=1$. The compiler and calibrated machine theorem establish
\begin{equation}\label{eq:source}
 \exists\mathbf n\in\N^{31}\ \bigwedge_{j=1}^{22}q_j(\mathbf n)=0
 \quad\longleftrightarrow\quad
 \text{a finite ZFC refutation exists}.
\end{equation}
The archive includes the exact sparse equations and their coordinate adapters.
The older transport coordinate is $z_{\rm old}=z+1$; this is a genuine positive
zero-set bijection, justified by the inherited compiler bounds, not by an
unrestricted substitution alone.

\subsection{A complete coefficient instance, not four free parameters}

The active four-port instance has radix $B_0=2^{1064921}$. Its port bit lengths
are
\[
 (1064921,\ 1070004,\ 1064921,\ 1064922).
\]
All four actual unsigned little-endian payloads are included, with SHA-256
hashes, and must be used together. Their raw total is 533,099 bytes.
The scalar source coefficients are explicit short integer polynomials in these
ports; every fixed coefficient leaf is therefore specified without referring
to halting or consistency. The new positional source does not use an unspecified
recursively enumerable index. Its additional large leaves are explicit powers
of two with fixed exponents, not exponentials in unknowns.

The mathematical dependency chain is
\[
\begin{split}
\text{new polynomial has a zero}
&\longleftrightarrow\text{all shifted quadratic rows have a zero}\\
&\longleftrightarrow\text{literal K215 halts on blank tape}\\
&\longleftrightarrow\text{ZFC has a finite refutation}.
\end{split}
\]
The new source adapter and norm modifications are proved here. The native
compiler, the established positive-system endpoint, and the relevant published
Pell/Lucas bridge theorems remain identified dependencies. Finite circuit checks
are not substitutes for these unbounded mathematical arguments. We did not
compile a new end-to-end Lean or Rocq development, and we do not infer an internal
ZFC theorem merely from the general soundness of either proof assistant.

\section{A 31-position certificate of maximum 799}

The positions, in source-variable order, are
\begin{equation}\label{eq:weights}
\begin{split}
(&10,35,55,62,105,144,145,159,168,232,245,306,323,325,353,371,\\
 &374,382,494,499,525,566,578,610,647,681,685,741,777,783,799).
\end{split}
\end{equation}
Append $w_0=0$. The 528 unordered pairs with repetition from these 32
positions have distinct sums modulo 1023. Thus they also have distinct sums
as integers. Padding a degree-zero or degree-one monomial with zero shows
that all monomials of total degree at most two have different weighted
exponents $\alpha\cdot\mathbf w$.

This property is an exceptionally small finite certificate: enumerate
$0\le i\le j\le31$, compute $w_i+w_j\pmod{1023}$, and reject repetitions.
The downstream proof needs only this explicit property, not a trusted
finite-field implementation or an optimization program.

For provenance, the positions were obtained by the Bose--Chowla method in
$\mathbb F_{32^2}$, represented by $\mathbb F_2[t]/(t^{10}+t^3+1)$.
A primitive element is $t$ and the subfield of 32 elements is generated by
$t^{33}$. Taking logarithms of $t+a$ over that subfield gives a modular
$B_2$ set. An affine unit transformation, multiplier 47 and translation
$-879$ modulo 1023, followed by ordering yields the displayed list with zero.
No optimality is claimed: the direct 528-pair verification is the correctness argument.

\begin{lemma}[Why the classical construction works]
If $\theta$ has degree two over $\mathbb F_q$, the logarithms of
$\theta+a$, for $a\in\mathbb F_q$, form a $B_2$ set in the multiplicative
cyclic group of $\mathbb F_{q^2}$.
\end{lemma}
\begin{proof}
Equality of two logarithm sums gives
$(\theta+a)(\theta+b)=(\theta+c)(\theta+d)$.
Subtracting the two monic quadratics leaves a polynomial of degree at most
one vanishing at $\theta$. It must be zero, and the multisets of roots agree.
Affine unit maps preserve equality of equal-length modular sums.
\end{proof}

\section{The multi-output coefficient theorem}\label{sec:multi}

Put $W=799$ and $h=2W+1=1599$. Reorder the residuals as
\[
1,2,3,4,6,7,\ldots,22,5,
\]
so the final row has constant coefficient one. For the row in position $j$,
assign target $H_j=jh$, where $1\le j\le22$. Hence
\[
 H_*=35178,\qquad K_*=H_*+2W=36776.
\]
Write that row as $q_j(\mathbf n)=\sum_{|\alpha|\le2}p_{j,\alpha}\mathbf n^\alpha$
(the reordered row index is understood), and define
\begin{align}
 c_{j,\alpha}&=(2-|\alpha|)!\prod_i\alpha_i!\,p_{j,\alpha},\label{eq:c}\tag{C}\\
 D_j(t)&=\sum_{|\alpha|\le2}c_{j,\alpha}t^{H_j-\alpha\cdot\mathbf w},
 &D_*(t)&=\sum_{j=1}^{22}D_j(t).\label{eq:D}
\end{align}
These are ordinary integer polynomials with fixed exponents. The archive's
136 records give the complete nonzero coefficient list.

\begin{theorem}[Simultaneous extraction]\label{thm:extract}
For arbitrary integers $n_1,\ldots,n_{31}$, let
$C(t)=(1+\sum_i n_i t^{w_i})^2$. Then
\[
 [t^{H_j}]\,C(t)D_*(t)=2q_j(\mathbf n)\qquad(1\le j\le22).
\]
If $M\ge1$ bounds the absolute values of all $c_{j,\alpha}$, every
coefficient of $D_*$ has magnitude at most $M$. If $0\le n_i<b$, every
coefficient of $CD_*$ has magnitude at most $M(31b)^2$.
\end{theorem}
\begin{proof}
The support of $D_j$ lies in $[H_j-2W,H_j]$; these intervals are disjoint.
The pair-sum certificate makes all positions within each interval distinct.
This proves the coefficient bound for $D_*$ without a factor of 22.
A contribution from row $k$ to target $H_j$ requires
\[
 H_j-H_k=\beta\cdot\mathbf w-\alpha\cdot\mathbf w,
 \qquad |\alpha|,|\beta|\le2.
\]
The right side belongs to $[-2W,2W]$, while every nonzero difference on the
left has magnitude at least $2W+1$. Thus $j=k$. The pair-sum property then
forces $\alpha=\beta$, and the multinomial coefficient cancels the factorial
multiplier in \eqref{eq:c}, leaving exactly $2q_j$.
For a fixed coefficient of $CD_*$, the triangle inequality bounds it by
$M$ times the sum of the nonnegative coefficients of $C$:
\[
 M\bigl(1+\textstyle\sum_i n_i\bigr)^2\le M(31b)^2.
\]
Off-target products of different rows may overlap. The last bound already
accounts for them; it does not incorrectly assume all product blocks disjoint.
\end{proof}

\begin{lemma}[Positivity before decoding]\label{lem:Dpositive}
For every integer $B>M+1$, $D_*(B)>0$.
\end{lemma}
\begin{proof}
The leading coefficient at degree $H_*$ is two, because the final shifted
row has constant one. Every other exponent is smaller. Therefore
\[
 D_*(B)\ge2B^{H_*}-M\frac{B^{H_*}-1}{B-1}>B^{H_*}>0.
\]
\end{proof}
This unconditional positivity is important for the simpler nine-variable
majorants. It does not depend on the encoded equations already being true.

\section{One carry test for all 22 equations}\label{sec:carry}

Let $b\ge4$ and let $B$ be a power of two satisfying
\[
 B\ge4M32^2b^2.
\]
Set
\[\begin{gathered}
\ell=\sum_iB^{w_i},\qquad \mathcal T=\sum_{j=1}^{22}B^{H_j},\\
N_0=B^{W+1},\qquad N_1=4B^{K_*+1}\\
N=N_0N_1,\\
M_0=N_0-1-(b-1)\ell,\\
T_0=M_0+(B-2)\mathcal T N_0,\\
J_0=2(1+g)^2D_*(B)+B\sum_{e=0}^{K_*}B^e,\\
S_0=g+J_0N_0,\\
\mathcal R=(N-1)\bigl((S_0+T_0+1)N+T_0+1\bigr).
\end{gathered}\]
Every sum limit and exponent is fixed. The arithmetic circuits evaluate
geometric sums by division-free doubling identities.

Suppose $0\le g<2bB^W$. The bounds on $B$ imply $g<N_0$ and
$1+g\le3bB^W$. By Lemma~\ref{lem:Dpositive}, $J_0>0$; the coefficient
bound for $D_*$ gives
\[
 J_0<\frac{3B}{2}\frac{B^{K_*+1}}{B-1}<N_1.
\]
Also $0\le M_0<N_0$ and $(B-2)\mathcal T<N_1$. Thus
$0\le S_0,T_0<N$. If $b$ is a power of two, the mask $M_0$ has base-$B$
digits $B-b$ in the 31 allowed cells and $B-1$ elsewhere.

\begin{lemma}[Binary-carry identity]\label{lem:carry}
For $N=2^r$ and $0\le S,T<N$, put
$X=(N-1)((S+T+1)N+T+1)$. Then
\[
 N^2\mid\binom{2X}{X}\quad\longleftrightarrow\quad
 \text{$S+T$ has no binary carry}.
\]
\end{lemma}
\begin{proof}
The exact identity
\[
 X=(S+T)N^2+(N-1-S)N+(N-1-T)
\]
has disjoint lower blocks, even when $S+T\ge N$. Therefore
\[
 \pc(X)=\pc(S+T)+2r-\pc(S)-\pc(T)=2r-\tau_2(S,T),
\]
where $\tau_2$ counts binary carries. Legendre's formula gives
$v_2\binom{2X}{X}=\pc(X)$, proving the assertion.
\end{proof}

\begin{theorem}[Multi-target source interface]\label{thm:sourcecarry}
Under the preceding hypotheses, with $b$ a power of two,
\[
 N^2\mid\binom{2\mathcal R}{\mathcal R}
 \quad\longleftrightarrow\quad
 \exists\mathbf n\in[0,b)^{31}\left[
 g=\sum_i n_iB^{w_i}\ \land\ \bigwedge_jq_j(\mathbf n)=0\right].
\]
\end{theorem}
\begin{proof}
Lemma~\ref{lem:carry} splits the no-carry test into the low block
$(g,M_0)$ and the high block $(J_0,(B-2)\mathcal T)$. The first test decodes
exactly the displayed 31 digits. This does not assume that the low-block sum
has no carry into its next block: any such carry already fails the test.

For the decoded digits, put $a_e=[t^e](CD_*)+B/2$. The coefficient bound
makes $0<a_e<B$, and
\[
 J_0=2\sum_{e=0}^{K_*}a_eB^e.
\]
Divide both high-block summands by two, merely shifting their binary bits.
The test becomes carry-freeness of $\sum a_eB^e$ and
$(B/2-1)\mathcal T$. At a target cell, the second summand has every low
binary bit set. Because $0<a_{H_j}<B$, carry-freeness is equivalent to
$a_{H_j}=B/2$. By Theorem~\ref{thm:extract}, this means $q_j(\mathbf n)=0$.
Non-target cells have zero mask and impose no extra equation.
\end{proof}

Every simultaneous nonnegative zero is nonzero because $q_5(0)=1$.
Choosing $b$ as an arbitrarily large suitable power of two gives
$g\ge B\ge b$ and $g<bB^W$. Thus the source interface has both soundness
and the cofinal completeness needed by the arithmetic bridges. No unknown
index, extra source witness, or variable exponent is introduced.

\paragraph{A source-transport subtlety.} The optimized parent uses $z_{\rm old}=z+1$ before its positive coordinates are shifted to nonnegative ones. This is not a bijection on positive witnesses. Soundness is immediate, but completeness uses the canonical compiler bound $z_{\rm old}\ge2$, proved in the pinned \texttt{quartic31/QUARTIC60\_PROOF.md}. That proof is an explicit dependency of both new source specializations.

\section{The nine-variable equation}
\subsection{A fully polynomial specification}
Use the mixed domain $\epsilon,\phi,n\in\N$ and $g,h,s,w,i,j\in\Pos$.
Let $b=2^{2134939}+\epsilon$ and $B=b^3$, and use the preceding source
registers $R,N$. The 136 signed coefficient recipes have absolute values
strictly below $M=2^{2134927}$, as follows by evaluating their absolute-value
bounds at the four port bit-length bounds. Thus $b\ge4096M$.
Put $G_0=32bB^{799}$ and define the following straight-line registers:
\[\begin{gathered}
U=N^2w,\qquad Y=N^2s\\
M_0=RY,\qquad P=2M_0^2U,\\
K=R+1+h(P-1),\qquad A=M_0(U+1)\\
B_1=2R+1,\qquad W_1=bw,\\
C=B_1+W_1+\phi,\qquad J=4A-5\\
Q=3W_1C-2(W_1^2-1),\\
D=(A^2-1)C^2+1,\qquad E=iJDC^2\\
F=(A^2-1)E^2+1,\\
G=A+F(CD+1-A),\qquad H=B_1+jC\\
I=(G^2-1)H^2+1.
\end{gathered}\]
(The letters here are local registers, not the fixed source ports.) Set
\[\begin{gathered}
a_0=(P^2-1)K^2+1,\qquad a_1=DFI\\
u=JF,\\
v=QF+J(H-C),\qquad m=K^2G_0-4(C-KY)^2G_0-8gK^2,\\
V_0=PK,\qquad V_1=(AC)(AE)(GH)\\
M_d=Q+H,\\
\Omega=um,\qquad \Theta=u(n+M_d)+v-\Omega(2M_d+V_0V_1).
\end{gathered}\]
Our equation is
\begin{equation}\label{eq:nine}
 \boxed{\ F_9=(\Theta^2-\Omega^2(a_0+V_0^2a_1))^2
                 -4\Omega^4V_0^2a_0a_1=0.\ }
\end{equation}
All displayed powers have fixed nonnegative integer exponents; the binomial
coefficient used in the proof is \emph{not} part of this polynomial.
For the nine-positive form replace $\epsilon,\phi,n$ by the corresponding
positive coordinates minus one. For the nine-nonnegative form instead replace
$g,h,s,w,i,j$ by their coordinates plus one. These are domain bijections,
not permission to interpret the same equation over unrestricted integers.

\subsection{The arithmetic bridge and the new finalizer}
The numerical ratio theorems in the pinned \texttt{NineRatio.lean} give the
following interface. For $R,N\ge8$, $0<b\le R,N$, $G_0>0$, and the indicated
witness domains, if $a_0,a_1$ are squares, $u\mid v$, and $m>0$, then
\[
 N^2\mid\binom{2R}{R},\qquad b\text{ is a power of two},\qquad8g<G_0.
\]
Conversely these conditions, with the stronger $32g<G_0$, supply the
arithmetic witnesses. This is the existing \emph{numerical} bridge, not its
old exponentially spaced source specialization. The complete source and
ratio proofs are included as pinned dependencies.

For the present source, positivity holds for every outer assignment:
$D(B)>0$, the low mask is positive, and $R\ge N\ge8$, $b\le R,N$.
On the soundness range $g<4bB^{799}$, the coefficient bound above gives
$S,T<N$, so the carry theorem applies. On the completeness range the source
encoding satisfies $g<bB^{799}$, which gives $32g<G_0$. This verifies both
uses of the bridge without assuming a bound that is only proved afterwards.

\begin{lemma}[Two-radical finalizer]
In the displayed domain, \eqref{eq:nine} has a witness $n\in\N$ if and only
if $a_0,a_1$ are squares, $u\mid v$, and $m>0$.
\end{lemma}
\begin{proof}
All $a_0,a_1,u,V_0,V_1$ are positive. Direct subtraction gives
$V_0^2-a_0=K^2-1>0$. The analogous three identities for $AC,AE,GH$ imply
$V_1^2>a_1$. Also
\[
 0<v/u=Q/J+(H-C)/F<Q+H-1=M_d-1.
\]
If $m=0$, then $\Theta=u(n+M_d)+v>0$, excluding a zero. Otherwise the norm
factorization gives $\Theta/\Omega=r_0+V_0r_1$ for real roots $r_t^2=a_t$.
If this rational sum is nonzero, squaring yields
$2(\Theta/\Omega)r_0=(\Theta/\Omega)^2+a_0-V_0^2a_1$; hence both roots are
rational, and then integral. A zero sum would imply
$|r_1|=|r_0|/V_0<1$, impossible for the square root of the positive integer
$a_1$. Thus the roots are integral and $|r_t|\le V_t-1$.
Consequently $Z=2M_d+V_0V_1+r_0+V_0r_1\ge2M_d+1$. The vanishing factor says
\[
 n+M_d+v/u=mZ.
\]
Its left side is positive, proving $m>0$; integrality proves $u\mid v$.
Conversely choose integral roots and set $n=mZ-M_d-v/u$. The same bounds give
$n\ge2$, and the appropriate factor is zero.
\end{proof}
Combining this lemma, the numerical bridge, and the simultaneous source
compiler proves the nine-variable equivalence in Theorem~\ref{thm:main}.

\subsection{Exact degree despite a leading cancellation}
Let $t=\deg N=3(H_{22}+3W+2)=112731$. On the diagonal substitution assigning
one indeterminate to every variable, $N$ has leading coefficient $4$ and
$R$ leading coefficient $16$ at degree $3t$. The important cancellation is
resolved by the exact positive identity
\[\begin{gathered}
V_1^2-a_1={}(C^2-1)(AE)^2(GH)^2\\
+D(E^2-1)(GH)^2+DF(H^2-1).
\end{gathered}\]
The first two summands have degree $298t+64$ and each has leading coefficient
$2^{518}$. Their sum has coefficient $2^{519}$; the third has lower degree.
Writing $E_0=u(n+M_d)+v-2\Omega M_d$ gives
\[
 \Theta^2-\Omega^2(a_0+V_0^2a_1)
 =E_0^2-2E_0\Omega V_0V_1+\Omega^2V_0^2(V_1^2-a_1)-\Omega^2a_0.
\]
The third term uniquely dominates, with degree $576t+6W+140$ and coefficient
$2^{1025}$. Hence $\deg F_9=1152t+12W+280=129875980$, with diagonal leading
coefficient $2^{2050}$. The circuit certificate also supplies the matching
multivariate upper bound. This is an exact degree, not merely a syntactic
estimate or a numerical experiment.

\section{Ten unrestricted integer variables}
\subsection{Source adaptation and equation}
Use integer variables $f,g,h,k,l,w,x,y,z,m$. Here
$b=1+3f$, $B=\beta_s b^2$, where the exact power-of-two constant $\beta_s$
is specified in \texttt{sun/MULTITARGET\_INSTANCE.json}. Use the same
simultaneous source registers and put $X=R$, $Y=N^2$, $C_0=bB^{799}$.
The following registers specify the unchanged Sun arithmetic tail:
\[\begin{gathered}
L=lY,\qquad U=2LX\\
V=4gwY,\qquad W_0=bw,\\
A=U(V+1),\qquad B_0=2X+1\\
C=B_0+(A-2)h,\\
P=U^2V,\qquad K=X+1+k(P-2)\\
S_0=2A-5,\\
T_0=3W_0C-2(W_0^2-1),\qquad r_1=(P^2-4)K^2+4,\\
\delta=8C_0^3gK^2-g^2[32(C-KL)^2C_0^3+g^2K^2],\\
O=f^2l^2\delta,\qquad \xi=zx\\
D=(A^2-4)C^2+4,\\
E=C^2D\xi,\qquad F=4(A^2-4)E^2+1,\\
G_*=2(A+2)(A-2)^2E^2-1,\qquad \Omega_*=16E,\\
H_*=C+(B_0-C)F+\Omega_*Fy,\qquad I_*=(G_*^2-1)H_*^2+1.
\end{gathered}\]
Set
\[\begin{gathered}
\Sigma=\xi S_0^5,\qquad r=T_0+S_0y\\
\vartheta=(r^3-rS_0^2-S_0^3)(r^2+23S_0^2),\\
r_0=b[1+(3O-4)b^2z^2],\qquad r_2=DFI_*,\\
w_0=1,\qquad w_1=r_0\\
w_2=r_0(P^2+K^2+3),\\
Z=\Sigma m-\vartheta,\qquad A_i=\Sigma^2w_i^2r_i\quad(0\le i\le2),\\
V_*=Z^2+A_2-A_0-A_1,\qquad W_*=V_*^2+4Z^2A_2-4A_0A_1.
\end{gathered}\]
The ten-integer equation is
\begin{equation}\label{eq:ten}
 \boxed{\ F_{10}=W_*^2-16V_*^2Z^2A_2=0.\ }
\end{equation}
Its exact straight-line expansion has 936 live binary operations. The
multivariate degree upper bound $237526776$ is attained: evaluation of the
leading homogeneous circuit at the diagonal gives nonzero residues $54$,
$453$, and $982432$ modulo $101$, $1009$, and $1000003$, respectively.
One nonzero residue suffices; three provide independent replay diagnostics.

\subsection{Why the new source satisfies the old tail contract}
The signed tail is used as a proved arithmetic lemma, with its complete
pinned proof supplied in \texttt{sun/upstream/SIGNED10\_PROOF.md}, Sections
4--8. Its requirements on the source are precisely the following:
\begin{enumerate}
\item If $b$ is a positive square, $f\ne0$, and $0\le g<2C_0$, the source registers have the required positive bounds. If in addition
$b$ is a power of two, the binomial divisibility condition implies a common
source zero. The reverse arithmetic bridge supplies that additional condition;
it does not follow from squareness alone.
\item Every source zero has arbitrarily large square choices $b=1+3f$ and
an encoding $b\le g<C_0$ meeting the stronger completeness bounds.
\item $3\mid X$ and the numerical size inequalities used by the interpolation
lemma hold.
\end{enumerate}
For (1), the signed positivity gadget first forces $b$ to be a positive
square; because $b=1+3f$ and $f\ne0$, it forces $f>0$ and $b\ge4$.
The fixed $\beta_s$ exceeds $4M32^2$, so the same coefficient and block
estimates prove (1). For (2), choose arbitrarily large even-exponent powers
of two congruent to $1$ modulo three. For (3), $B\equiv1\pmod3$ and
$N\equiv1\pmod3$, so the factor $N-1$ in $R$ gives $3\mid X$.
The complete bounds, including the unsigned and signed source ranges, are
proved in \texttt{sun/MULTITARGET\_SOURCE\_PROOF.md}.

For clarity, the tail is not a black-box appeal to an unspecified universal
polynomial. Its main ingredients are the following explicit facts.
The polynomial $\pi(t)=(t^3-t-1)(t^2+23)$ has a root modulo every positive
integer but has no rational root. The factor $S_0^5$ and a valuation argument
turn $\Sigma\mid\vartheta$ into $S_0\mid T_0$. Full-residue interpolation
supplies the converse witness in the required residue class. The square
condition on $r_0$, together with $z\ne0$, forces positivity and the square
condition on $b$; conversely Pell saturation provides $z$.
Finally \eqref{eq:ten} is the norm of
$Z+\Sigma\sum_{i=0}^2\varepsilon_iw_i\sqrt{r_i}$ over all eight sign choices.
The weights strictly separate the largest radical whose sign changes under
an \emph{actual} automorphism of the multiquadratic field. This forces each
radicand to be a square; it does not assume three independent sign
changes when the square classes are dependent. If $\Sigma=0$, the norm is
$\vartheta^8\ne0$, excluding the degenerate branch. The supplied proof then
recovers the numerical binomial condition and the source zero. Conversely
its interpolation and Pell lemmas produce every required witness.

The new result is the verified replacement of the source in this tail.
The general interpolation theorem is imported from Sun's 2026 preprint,
not newly reproved or machine-checked here. The older Jones-based nine-positive
construction is therefore the more conservative primary recommendation.

\section{A smaller low-degree twelve-positive equation}
This independent improvement is useful when coefficient extraction and
large degree are undesirable. It reduces the radical-separation weight rather
than the number of quantified coordinates. Let all twelve variables
$J,F,\alpha,z,h,i,T,s,w,\eta,\zeta,t$ be strictly positive and set
\[\begin{gathered}
q=\beta J+1,\qquad C=q-F-\alpha\\
X=wq,\qquad Y=sq^3,\\
E=XY,\qquad k=\eta+\zeta\\
c=kY+\eta,\qquad a=E+Y,\\
H=4a+3,\qquad \Delta=a^2+H\\
\ell=EkY,\\
A=1+\ell(\ell+k),\qquad G=1+\Delta c^2\\
b_0=X+ac,\\
R_0=((q-1)(q-F)+\alpha)(q^2-1)+(M_C+qM_F)J,\\
N_k=k-hE-R_0,\qquad N_t=(K_0+w)C+q-F-z(q-1),\\
I=ic^2,\qquad S=\Delta I\\
D=1+IS,\qquad Q=S^2\\
d=Q-1,\\
v=cT,\qquad b_1=c+R_0D\\
P=QDv^2,\qquad U=Qb_1^2,\\
u=k^4i,\qquad M=u(u+T)\\
L=HM+1,\qquad f_0=(a+2)I.
\end{gathered}\]
For the degree-941 form take
\[
 \sigma=b_1+vf_0+M((a+2)c-\ell),\qquad W=b_0+H(t+\sigma).
\]
Use the polynomial registers
\[\begin{gathered}
n_0=(HM)^2A,\qquad d_n=dn_0\\
P_H=H^2P,\qquad U_H=H^2U,\\
e_H=P_H+U_H-H^2,\qquad j_0=4P_HU_H,\\
T_r=2d(W^2+L^2G),\qquad K_n=[d(W^2-L^2G)]^2,\\
A_a=2(d_n+e_H),\qquad c_0=(d_n-e_H)^2-j_0,\\
a_t=T_r-A_a,\qquad b_t=c_0-K_n\\
c_1=16d_nj_0,\\
v_r=a_t(a_tT_r+2b_t)-c_1,\qquad u_r=b_t^2-a_t^2K_n,\\
\mathcal N=u_r(u_r+T_rv_r)+K_nv_r^2.
\end{gathered}\]
The single equation and its alternative three-equation system are
\begin{equation}\label{eq:twelve}
 \boxed{\ N_kN_t(1+\mathcal N)-1=0\ },\qquad
 \boxed{\ N_k=1,\quad N_t=1,\quad\mathcal N=0\ }.
\end{equation}
These have respectively 121 operations and degree 941, or 116 shared
operations and maximum degree 932. Sharing is counted across system outputs;
subtracting an imagined fixed number of gates from the scalar count gives an
incorrect system count.

\subsection{The uniform new separation estimate}
The inherited compiler interface gives, even on the weak branch
$N_k=N_t=\varepsilon\in\{1,-1\}$,
\[
 q\ge16,\quad X\ge q,\quad Y\ge q^3,\quad E=XY,\quad
 3q<R_0<q^4-2,\quad k=R_0+\varepsilon+hE>E.
\]
It also gives $kY<c<k(Y+1)$ and
$\sqrt G-\sqrt A>kY^2$, $\sqrt A>kY^2$. Consequently
\[
 Y<k/16,\qquad a+2<9k/8,\qquad c<17kY/16.
\]
The definitions yield the sharper estimates
\[
 \frac{v\sqrt D}{kY^2}<\frac{uT}{8},\qquad
 \frac{b_1}{kY^2}<1+\frac{u^2}{128}.
\]
Thus
\[
 2(v\sqrt D+b_1)<MkY^2
 <\min\{M\sqrt A,\ M(\sqrt G-\sqrt A)\}.
\]
This directly proves the radical separation required by the inherited
norm argument. It is \emph{not} necessary for the new weight to exceed the
old, much larger weight. The unchanged recovery argument forces the relevant
square roots to be rational integers, while the Pell saturation argument
supplies the converse witnesses. Modulo four the norm is a square, excluding
the otherwise possible $\mathcal N=-2$ branch of the product-unit equation.
For the shift, write the old witness as
$t_{\rm old}=\gamma+Mg+M\tau+y$. The existing inequalities give
$M\tau>vf_0+b_1$ and $(a+2)c-\ell<g$; hence
$t=t_{\rm old}-\sigma>0$. This verifies completeness for the shifted equation,
not just preservation of soundness.

The full norm recovery, Pell dependencies, and uniform estimate are provided
in the supplied positive-variable proof and its pinned dependencies.
The exact degree certificate resolves the leading cancellations via the eight
radical factors. Four have degree 119; the other four have degrees
117,117,115,107. Their sum is 932, and the two outer residual factors add
$7+2$, giving degree 941.

\begin{center}
\begin{tabular}{lrrrr}\toprule
Shift $\sigma$ & Scalar ops & Degree & System ops & Max. degree\\\midrule
$b_1$ &112&961&107&952\\
$b_1+vf_0$ &116&949&111&940\\
$b_1+vf_0+M(ac-\ell)$ &119&945&114&936\\
$b_1+vf_0+M((a+2)c-\ell)$ &121&941&116&932\\\bottomrule
\end{tabular}
\end{center}
No row dominates all the others. The twelve-variable proposal improves the
previous same-family degree/operation tradeoff but does not beat the repository's
67-operation scalar or 59-operation system when variable count and degree
are ignored.

\section{What is checked, what is imported, and what is not claimed}
\subsection{Constants are actual integers, not extra variables}
The four fixed compiler ports are supplied as compressed unsigned
little-endian byte strings, with raw-byte lengths and SHA-256 hashes in both
instance manifests. The same coherent package-31 tuple is used throughout;
its total raw size is 533099 bytes. Replacing one of these ports by a smaller
number from another package would not preserve the theorem.
Every coefficient recipe is a finite integer polynomial in those ports,
recorded exactly, not a semantic instruction to find a ZFC proof. All 17
live fixed leaves of each primary nine- and ten-variable DAG are additionally
materialized. Each has a sign byte followed by an unsigned little-endian
magnitude, compressed by gzip, and both raw and compressed hashes are given.
The nine-variable leaves occupy 3464120 raw bytes and 64286 compressed bytes;
the ten-variable leaves occupy 64286 compressed bytes. Duplicate values and
fixed integers count as coefficients, not quantified unknowns.

The fixed Sun radix is $\beta_s=2^{2134940}$; the Jones floor is
$2^{2134939}$. These are different constants serving different source
interfaces. Neither is the original fixed compiler port $\beta$.
The coefficient bounds are derived from the actual 136 recipes and the four
verified bit-length bounds, not from a hoped-for smaller tuple.

\subsection{Operation and degree conventions}
An operation is a live binary addition, subtraction, or multiplication in a
straight-line DAG, including multiplication by a fixed numeral and a
squaring. Fixed powers are expanded by multiplication chains. Dead gates
are excluded; shared registers are counted once, including across system
outputs. Literal and fixed-coefficient leaves are not arithmetic operations.
Consequently operation counts do not include the offline cost of generating
the fixed coefficients, nor the bit complexity of evaluating giant integers.
The equations are succinct polynomial specifications, not dense monomial
expansions. The exact DAG, rather than typography in this report, is the
canonical machine-readable polynomial definition.

The Sun source is evaluated by homogenizing each quadratic row with
$B^{W-w_i}$ and then Horner-combining the target rows. The identity
\[
 \sum_j B^{H_j-2W}\!\sum_{|\alpha|\le2}c_{j,\alpha}
 (B^W)^{2-|\alpha|}\prod_i(B^{W-w_i})^{\alpha_i}
 =\sum_{j,\alpha}c_{j,\alpha}B^{H_j-\alpha\cdot w}
\]
proves that this is exactly the same $D(B)$, not a changed equation.
It cuts a preliminary 954-operation Sun DAG to the delivered 936-operation
DAG. The nine-variable evaluator uses the same homogeneous row strategy.

\subsection{The proof dependency chain}
The logical chain is explicit:
\[
 \begin{gathered}
 \text{finite ZFC refutation}\ \Longleftrightarrow\
 \text{fixed compiler source zero}\\
 \Longleftrightarrow\ \text{simultaneous quadratic carry witness}\\
 \Longleftrightarrow\ \text{numerical Jones or Sun bridge witness}\\
 \Longleftrightarrow\ F_9=0\text{ or }F_{10}=0.
 \end{gathered}
\]
The new middle source construction and its transport conditions are proved
here and in the detailed proof files. The first equivalence depends on the
repository's fixed K215/compiler endpoint and coefficient binding; the
arithmetic bridge depends on the identified Jones results, or on Sun's
2026 interpolation theorem for the ten-integer construction. The final norm
arguments are explicit. These are ordinary mathematical proofs using
arithmetical arguments formalizable in ZFC. They do not establish
$\Con(\mathrm{ZFC})$ itself, and do not assume it in either direction.

No new complete Lean/Rocq certificate or checked object-level derivation
inside a formal ZFC proof system is delivered. Existing checked theorems do
not automatically certify the new coefficient specialization or source
replacement. The finite verifiers below check exact polynomial identities,
source transcription, and certificates; they cannot replace the universal
proof or determine whether the final equations have a zero.

\subsection{Reproducible audit inventory}
The archive contains the authored finite verifiers, exact JSON DAGs, source
pins, manifests, fixed coefficient payloads, and inherited proof texts.
The most important audits are:
\begin{itemize}
\item The independent quadratic checker reconstructs all 136 weighted
coefficient records from the 22 original residuals; verifies all 528 pair
sums; and checks 71808 cross-target incidences.
\item The nine-variable checker verifies domain translations, 180
prime-field circuit/certificate replays, 15 actual-port replays, 128
finalizer completeness fixtures, 980 zero-margin exclusions, 352 independent
extractions, all 17 fixed payloads, and the exact symbolic degree leader.
\item The ten-variable checker reconstructs its coefficients independently,
replays the homogeneous source identity, checks its payload hashes and actual
port binding, and verifies nonzero leading residues at the exact degree.
\item The positive-variable checks compare 240 modular radical products and
1080 explicit upper-bound cases; a separate review checks the uniform
inequality argument and the source domains.
\item The final comparison inventory authenticates 84 pinned source files by
Git blob identifiers. The authored short-form audit records failed local
simplification ideas rather than presenting them as improvements.
\end{itemize}
The final release includes a hash manifest and a driver for running the
checks in an isolated temporary copy. Only newly authored scripts are run;
upstream scripts, where retained for provenance, are inert source material.
See \texttt{README.md} for exact commands and required Python packages.
The JSON check outputs record actual results, while this report states the
universal proof obligations separately. No measured build-speed improvement,
global minimum, or unbounded computational search result is claimed.

\section{Research directions that remain genuinely open}
\begin{enumerate}
\item \textbf{Optimize the target layout jointly with the weights.}
The equal gap $2W+1$ is a simple sufficient condition. A finite certificate
could allow selected overlapping intervals when no cross-row monomial reaches
another target. Such a certificate must account for \emph{every} source
monomial and signed coefficient; random success is not a proof.
\item \textbf{Remove source registers before positional encoding.}
An elimination that preserves quadracity and the canonical completeness
witness might reduce both the number of positions and the number of target
cells. Substitution that raises degree can erase the gain. The objective
should be the resulting bridge degree, not just parent variable count.
\item \textbf{Specialize the coherent tuple only after changing the source.}
The package-36 tuple may reduce coefficient storage if its complete compiler
binding can be transported to this exact parent. Mixing individual constants
is invalid. A separate full rebinding audit is required.
\item \textbf{Formalize the source interface once.}
A generic theorem for finitely many quadratic outputs, carrying an explicit
pair-sum certificate and target-isolation certificate, could replace large
instance-specific positional proofs in both Lean and Rocq. The new finite
checks suggest a compact certificate boundary but are not themselves that
formalization.
\item \textbf{Seek a simpler arithmetic bridge.}
Most remaining degree growth comes after the carry test. Removing one Pell
or norm layer could matter more than further improving the positions.
Any such proposal must retain unrestricted-integer soundness and exclude
zero multipliers; an improvement only on positive witnesses does not suffice
for the ten-integer theorem.
\end{enumerate}
The direct multi-output construction is a structural improvement, not merely
a smaller coefficient or a renamed polynomial. It preserves the established
variable counts while replacing exponentially large positional degrees by
explicit moderate integers and reducing their live evaluation circuits.

\section*{Postscript: a newer formalization-status commit}
After the construction snapshot, commit
\texttt{8747904f8b973f9016434993e7ec6ae6ec7549ef} was reviewed on
23 October 2026. It adds \texttt{Computability/DiophantineConsistency/README.md}
and \texttt{STATUS.md}, distinguishing mathematical equivalences, internal
portability, and checked object-level ZFC derivations. It adds no new equation
or checked Lean theorem. The construction pin remains unchanged. Its caution
is adopted here: these new equivalences are mathematical results with explicit
dependencies, not newly delivered kernel-checked derivations.

\begin{thebibliography}{9}
\bibitem{repo} V. Reshetnikov, \emph{ProveIt}, pinned source \pin.
\url{https://github.com/VladimirReshetnikov/ProveIt}.
Exact file paths and blob identifiers are in \texttt{review/SOURCE\_PINS.json}.
\bibitem{jones} J. P. Jones, Universal Diophantine equation,
\emph{Journal of Symbolic Logic} \textbf{47} (1982), 549--571.
\url{https://doi.org/10.2307/2273588}.
\bibitem{bose} R. C. Bose and S. Chowla, Theorems in the additive theory of
numbers, \emph{Commentarii Mathematici Helvetici} \textbf{37} (1962/63).
\url{https://repository.ias.ac.in/8578/}.
\bibitem{sun21} Z.-W. Sun, Further results on Hilbert's tenth problem,
\emph{Science China Mathematics} \textbf{64} (2021), 281--306;
preprint version 7. \url{https://arxiv.org/abs/1704.03504v7}.
\bibitem{sun26} Z.-W. Sun, \emph{Ten unknowns for Hilbert's tenth problem over the integers},
2026 preprint, version 2. \url{https://arxiv.org/abs/2609.01594v2}.
The exact interpolation dependency is identified in the supplied signed-tail
proof; no new verification of that full preprint is claimed.
\end{thebibliography}
\end{document}
